% =============================================================================
%  PAPER SCAFFOLD — GRPO reward-component ablation
%
%  HOW TO USE
%  - Upload this folder to Overleaf (New Project -> Upload Project) and compile
%    with pdfLaTeX. It compiles as-is with the standard TeX Live on Overleaf.
%  - Sections 1-4 (Intro, Related Work, Method, Setup) are DRAFTED from your repo
%    and writable now. Sections 5-6 (Results, Analysis) are SCAFFOLDED with
%    placeholder tables/numbers to fill once the sweep finishes.
%  - Find every gap by searching for \todo and \res in this file.
%  - Before submission: (1) swap \documentclass to your target venue's style
%    (NeurIPS/ICLR/ACL workshop), (2) VERIFY every \cite entry in references.bib
%    (arXiv IDs/years are from memory and must be checked), (3) add real figures.
% =============================================================================
\documentclass[11pt]{article}

\usepackage[utf8]{inputenc}
\usepackage[T1]{fontenc}
\usepackage[margin=1in]{geometry}
\usepackage{amsmath,amssymb,amsfonts}
\usepackage{booktabs}
\usepackage{graphicx}
\usepackage{xcolor}
\usepackage{multirow}
\usepackage{algorithm}
\usepackage{algpseudocode}
\usepackage[numbers]{natbib}
\usepackage{hyperref}
\hypersetup{colorlinks=true, linkcolor=blue, citecolor=blue, urlcolor=blue}

% --- authoring helpers: remove (or redefine to {}) before final submission ----
\newcommand{\todo}[1]{\textcolor{red}{[\textbf{TODO:} #1]}}
\newcommand{\res}[1]{\textcolor{blue}{\textit{#1}}}   % wrap placeholder numbers
% -----------------------------------------------------------------------------

\title{\textbf{Which Rewards Matter? A Controlled Component Ablation of\\
GRPO for Small-Model Mathematical Reasoning}}

\author{
  Moh. Izzatul Haqi\thanks{Department of Information Technology,
  Institut Teknologi Sepuluh Nopember (ITS), Surabaya, Indonesia.
  Correspondence: \texttt{<email>}.} \quad
  \todo{Advisor name (co-author)} \\
  \todo{affiliation / lab}
}
\date{}

\begin{document}
\maketitle

\begin{abstract}
Reinforcement learning with verifiable rewards (RLVR) has become a standard
recipe for eliciting mathematical reasoning in language models, and the reward
functions used in practice increasingly combine several signals---answer
correctness, output format, length control, and partial-credit terms. Yet the
\emph{marginal} contribution of each component is rarely isolated, and results
are frequently reported from a single training run, leaving open which signals
genuinely help and which merely invite reward hacking. We present a controlled,
multi-seed ablation of the reward function in Group Relative Policy Optimization
(GRPO) on a small (1.5B-parameter) model trained under a fixed QLoRA budget. We
decompose the reward into correctness, format, length, and partial-credit
components and evaluate a cumulative ladder of five variants (A0--A4) on GSM8K
and on out-of-distribution MATH-500, reporting paired significance tests and
effect sizes across seeds. We find that \todo{key finding 1 (e.g.\ which
component yields the largest pass@1 gain)}, that \todo{key finding 2 (e.g.\ a
component that helps the reward signal but not held-out accuracy)}, and that
\todo{key finding 3 (e.g.\ the pass@1 vs.\ pass@k pattern)}. Our results
\todo{one-sentence takeaway}. Code is available at
\url{https://github.com/mohizzatulhaqi/grpo-reward-ablation}.
\end{abstract}

% =============================================================================
\section{Introduction}
% --- WRITABLE NOW; no results needed ---
Reinforcement learning has emerged as the dominant approach for improving the
reasoning ability of large language models (LLMs). Where early work relied on
learned reward models trained from human preferences~\citep{instructgpt}, recent
systems for mathematical and programmatic reasoning instead use
\emph{verifiable} rewards: a rule checks whether the final answer is correct,
yielding a sparse but unambiguous training signal~\citep{deepseekr1}. Group
Relative Policy Optimization (GRPO)~\citep{deepseekmath} has made this recipe
particularly practical by estimating advantages from a group of sampled
completions, removing the need for a separate value network and substantially
lowering the cost of RL fine-tuning.

In practice, the reward used to train reasoning models is rarely a single term.
Implementations commonly add a \emph{format} reward to encourage parseable
output, a \emph{length} penalty to curb verbosity, and various
\emph{partial-credit} or process signals intended to densify the otherwise
sparse, end-of-sequence reward and improve credit
assignment~\citep{lightman2023}. Each addition is plausible in isolation, but
the components interact, and a poorly designed term can be exploited: a model
may learn to satisfy a proxy (e.g.\ produce a well-formatted or short answer)
without improving its actual accuracy. Despite the centrality of these design
choices, their individual contributions are seldom measured in a controlled way,
and the RLVR literature frequently reports single-seed results, making it
difficult to separate genuine effects from run-to-run variance.

This paper asks a deliberately narrow but practically important question:
\emph{within a fixed GRPO setup, which reward components actually improve
reasoning, which are redundant, and which induce reward hacking?} We answer it
with a cumulative ablation on a small (1.5B) model. Starting from a pure
correctness reward (RLVR), we add one component at a time---format, length, and
partial credit---defining five variants A0 through A4. Crucially, we hold the
optimization budget fixed (the same QLoRA configuration, training steps, and
group size across variants) so that differences are attributable to the reward,
and we run each variant across multiple random seeds, reporting paired
significance tests and effect sizes rather than point estimates.

We deliberately study the small-model, limited-compute regime. This setting is
both increasingly relevant---efficient post-training of small open models is a
practical concern---and methodologically convenient: a general-purpose (rather
than math-specialized) base model leaves headroom for the reward to move
accuracy, and an out-of-distribution evaluation on MATH-500 tests whether gains
generalize beyond the training distribution.

\paragraph{Contributions.}
\begin{itemize}
  \item We formulate a controlled, cumulative reward ablation (A0--A4) for GRPO
  and instantiate it with four reward components whose definitions are given in
  Section~\ref{sec:method}.
  \item We adopt a statistical protocol---multiple seeds with paired Wilcoxon
  signed-rank tests and Cliff's $\delta$ effect sizes---that is rarely applied
  in the RLVR setting, and we report a reward-hacking diagnostic that tracks each
  component's reward against held-out accuracy.
  \item Empirically, we find that \todo{summarize 2--3 findings once results are
  in}. All code and configurations are released for reproducibility.
\end{itemize}

% =============================================================================
\section{Related Work}
% --- WRITABLE NOW ---
\paragraph{RL for LLM reasoning.}
Reinforcement learning from human feedback established RL as a tool for aligning
LLM behaviour~\citep{instructgpt}, typically optimized with Proximal Policy
Optimization (PPO)~\citep{ppo}. For tasks with checkable answers, RL with
verifiable rewards replaces the learned reward model with a deterministic
correctness check, which DeepSeek-R1 used to elicit strong reasoning without
supervised reasoning traces~\citep{deepseekr1}. GRPO~\citep{deepseekmath}
simplifies the optimizer by computing group-relative advantages, and is now a
common choice for reasoning-oriented fine-tuning.

\paragraph{Reward design and credit assignment.}
Because the verifiable reward is sparse---granted only at the end of a long
generation---a recurring theme is densifying or shaping it. Process-supervised
rewards score intermediate steps and can outperform outcome-only supervision on
mathematical reasoning~\citep{lightman2023}. Format and length terms are widely
used in practice to stabilize training and control output, but they are proxies
whose effect on true accuracy is not guaranteed. Our partial-credit component is
a lightweight, rule-based instance of this densification idea, and our ablation
is precisely an attempt to quantify when such shaping helps.

\paragraph{Does RL expand or sharpen reasoning?}
A active debate concerns whether RLVR expands a model's reasoning capability or
merely concentrates probability on solutions the base model could already
sample. Evidence that RL improves pass@1 while leaving pass@$k$ (for large $k$)
largely unchanged suggests a sharpening rather than expansion
effect~\citep{yue2025}. We connect to this question directly by reporting
pass@$k$ alongside pass@1 for every variant.

\paragraph{Efficient post-training.}
Parameter-efficient methods make this study feasible on a single consumer GPU.
LoRA adapts a frozen model with low-rank updates~\citep{lora}, and
QLoRA~\citep{qlora} further quantizes the base model to 4-bit, enabling
RL fine-tuning of billion-parameter models in modest memory. We use QLoRA
uniformly across all variants; because the comparison is \emph{relative}, a
fixed parameter-efficient setup does not confound the ablation.

% =============================================================================
\section{Method}
\label{sec:method}
% --- WRITABLE NOW (from repo) ---

\subsection{Background: GRPO}
For a prompt $q$, GRPO samples a group of $G$ completions
$\{o_1,\dots,o_G\}$ from the current policy $\pi_{\theta_{\text{old}}}$ and scores
each with a scalar reward $r_i = R(o_i)$. Advantages are estimated by normalizing
rewards \emph{within} the group,
\begin{equation}
  \hat{A}_i = \frac{r_i - \operatorname{mean}(\{r_1,\dots,r_G\})}
                   {\operatorname{std}(\{r_1,\dots,r_G\}) + \varepsilon},
\end{equation}
which removes the need for a learned value function. The policy is then updated
with a clipped objective and a KL penalty toward a reference policy
$\pi_{\text{ref}}$:
\begin{equation}
  \mathcal{L}_{\text{GRPO}}(\theta) =
  -\,\mathbb{E}\!\left[
  \frac{1}{G}\sum_{i=1}^{G}\frac{1}{|o_i|}\sum_{t=1}^{|o_i|}
  \min\!\big(\rho_{i,t}\hat{A}_i,\;
  \operatorname{clip}(\rho_{i,t}, 1-\epsilon, 1+\epsilon)\hat{A}_i\big)
  - \beta\, \mathrm{D}_{\mathrm{KL}}\!\big(\pi_\theta \,\|\, \pi_{\text{ref}}\big)
  \right],
\end{equation}
where $\rho_{i,t} = \pi_\theta(o_{i,t}\mid q, o_{i,<t}) /
\pi_{\theta_{\text{old}}}(o_{i,t}\mid q, o_{i,<t})$ is the per-token importance
ratio.

\subsection{Reward components}
We decompose the reward into four components, each scaled to $[0,1]$ so that the
weights below act as clean relative importances. Let $\mathrm{ext}(o)$ denote the
extracted final answer of completion $o$, $a^\star$ the gold answer, and
$\mathcal{N}(\cdot)$ the multiset of numbers appearing in a string.

\begin{itemize}
  \item \textbf{Correctness} (the only term in A0):
  $r_{\text{corr}}(o) = \mathbb{1}\big[\mathrm{ext}(o) = a^\star\big].$
  \item \textbf{Format} (structure, independent of value):
  $r_{\text{fmt}}(o) = \tfrac12\,\mathbb{1}[\text{reasoning block present}]
  + \tfrac12\,\mathbb{1}[\text{parseable boxed answer present}].$
  \item \textbf{Length} (soft brevity, budget $T$, hard cut $H$):
  $r_{\text{len}}(o) = \operatorname{clip}\!\big(1 - \tfrac{\max(0,\,|o|-T)}{H-T},\,0,\,1\big).$
  \item \textbf{Partial credit} (process signal via intermediate-number overlap
  with the gold solution $g$):
  $r_{\text{part}}(o) = \dfrac{|\mathcal{N}(o) \cap \mathcal{N}(g)|}{|\mathcal{N}(g)|}.$
\end{itemize}

The full reward is a weighted sum $R(o) = \sum_{k} w_k\, r_k(o)$ over the active
components of a given variant.

\subsection{Ablation design}
The five variants form a cumulative ladder (Table~\ref{tab:variants}): each adds
one component to the previous. A0 is pure RLVR and serves as the internal
baseline against which all others are compared; A4 uses the full reward with
weights tuned on a single-seed sweep. Algorithm~\ref{alg:grpo} summarizes the
training loop. The format and correctness terms are intentionally orthogonal---a
completion can be well-formatted yet wrong---so that ``follows the format'' and
``gets the right answer'' can be measured separately.

\begin{table}[t]
  \centering
  \caption{Reward components active in each ablation variant. Weights are
  starting values tuned on a single-seed sweep; \res{report final A4 weights}.}
  \label{tab:variants}
  \begin{tabular}{lccccl}
    \toprule
    Variant & Correctness & Format & Length & Partial & Description \\
    \midrule
    A0 & \checkmark &            &            &            & pure RLVR baseline \\
    A1 & \checkmark & \checkmark &            &            & + format \\
    A2 & \checkmark & \checkmark & \checkmark &            & + length control \\
    A3 & \checkmark & \checkmark & \checkmark & \checkmark & + partial credit \\
    A4 & \checkmark & \checkmark & \checkmark & \checkmark & full, tuned weights \\
    \bottomrule
  \end{tabular}
\end{table}

\begin{algorithm}[t]
\caption{GRPO training with composite reward (one variant)}
\label{alg:grpo}
\begin{algorithmic}[1]
\Require base policy $\pi_\theta$, dataset $\mathcal{D}$, group size $G$,
         active components and weights $\{w_k\}$
\For{each training step}
  \State sample prompt $q \sim \mathcal{D}$
  \State sample group $\{o_1,\dots,o_G\} \sim \pi_{\theta_{\text{old}}}(\cdot\mid q)$
  \For{$i = 1$ to $G$}
    \State $r_i \gets \sum_k w_k\, r_k(o_i)$ \Comment{composite reward}
  \EndFor
  \State $\hat{A}_i \gets (r_i - \operatorname{mean}(r)) / (\operatorname{std}(r) + \varepsilon)$
  \State update $\theta$ by minimizing $\mathcal{L}_{\text{GRPO}}$
\EndFor
\end{algorithmic}
\end{algorithm}

% =============================================================================
\section{Experimental Setup}
\label{sec:setup}
% --- WRITABLE NOW (from repo) ---

\paragraph{Models.}
Our primary model is Qwen2.5-1.5B-Instruct~\citep{qwen25}; we deliberately use
the general instruction-tuned model rather than the math-specialized variant so
that the reward has room to change accuracy. Pipeline prototyping uses
Qwen2.5-0.5B-Instruct.

\paragraph{Data.}
We train on GSM8K~\citep{gsm8k} (7{,}473 grade-school math problems), whose
answers are trivially verifiable. We evaluate on the GSM8K test split
(1{,}319 problems) and, as an out-of-distribution test, on
MATH-500~\citep{lightman2023}, a 500-problem subset of the MATH
dataset~\citep{mathdataset}.

\paragraph{Training.}
All variants are trained with GRPO, implemented in the TRL
library~\citep{trl}, under an identical budget: 4-bit NF4
quantization with LoRA (rank \res{16}, $\alpha=\res{32}$) applied to all
attention and MLP projections~\citep{lora,qlora}, group size $G=\res{8}$, KL
coefficient $\beta=\res{0.04}$, learning rate \res{1e-6}, and \res{N} optimization
steps, with gradient checkpointing on a single 16\,GB GPU. The KL reference is
obtained by disabling the LoRA adapter, so no second full-size model is
instantiated.

\paragraph{Evaluation.}
For each test problem we sample $n=\res{8}$ completions at temperature
$\res{0.8}$. We report pass@1 (the unbiased estimator $c/n$ averaged over
problems) and pass@$n$ using the estimator of \citet{codex}, together with mean
completion length and the rate of format adherence.

\paragraph{Statistical protocol.}
Each variant is trained with \res{3} random seeds. We report mean $\pm$ standard
deviation and compare each variant to A0 with a paired Wilcoxon signed-rank test
(seeds matched across variants) and Cliff's $\delta$ effect size. We note that
with \res{3} seeds the smallest attainable paired $p$-value is bounded away from
conventional thresholds; we therefore foreground effect sizes and treat
$p$-values as supporting evidence rather than the primary claim.

\paragraph{Reproducibility.}
Reward functions, the verifier, the statistical analysis, and the resume-capable
training driver are released, with unit tests covering the components that
determine the reward signal and the reported metrics.

% =============================================================================
\section{Results}
\label{sec:results}
% --- SCAFFOLD: fill from results/ after the sweep ---
\todo{Replace the placeholder numbers in Tables~\ref{tab:main} and
\ref{tab:stats} with aggregated results, and add learning-curve and
accuracy--length figures.}

\begin{table}[t]
  \centering
  \caption{Main results (mean over seeds). \res{All numbers are placeholders.}}
  \label{tab:main}
  \begin{tabular}{lccccc}
    \toprule
    & \multicolumn{2}{c}{pass@1} & & & \\
    \cmidrule(lr){2-3}
    Variant & GSM8K & MATH-500 & pass@$n$ & mean len & format rate \\
    \midrule
    A0 & \res{--.-} & \res{--.-} & \res{--.-} & \res{---} & \res{--.-} \\
    A1 & \res{--.-} & \res{--.-} & \res{--.-} & \res{---} & \res{--.-} \\
    A2 & \res{--.-} & \res{--.-} & \res{--.-} & \res{---} & \res{--.-} \\
    A3 & \res{--.-} & \res{--.-} & \res{--.-} & \res{---} & \res{--.-} \\
    A4 & \res{--.-} & \res{--.-} & \res{--.-} & \res{---} & \res{--.-} \\
    \bottomrule
  \end{tabular}
\end{table}

\begin{table}[t]
  \centering
  \caption{Comparison to the A0 baseline on GSM8K pass@1: difference, paired
  Wilcoxon $p$, and Cliff's $\delta$. \res{Placeholders.}}
  \label{tab:stats}
  \begin{tabular}{lccc}
    \toprule
    Variant & $\Delta$ vs.\ A0 & $p$ (Wilcoxon) & Cliff's $\delta$ \\
    \midrule
    A1 & \res{+-.-} & \res{-.--} & \res{+-.--} \\
    A2 & \res{+-.-} & \res{-.--} & \res{+-.--} \\
    A3 & \res{+-.-} & \res{-.--} & \res{+-.--} \\
    A4 & \res{+-.-} & \res{-.--} & \res{+-.--} \\
    \bottomrule
  \end{tabular}
\end{table}

\paragraph{Per-component contribution.}
\todo{Report the cumulative ladder trend and (if run) a leave-one-out on A4 that
drops each component to measure its marginal effect, especially partial credit.}

\paragraph{Reward hacking.}
\todo{Present the diagnostic: for each variant, contrast the trajectory of each
component's reward with held-out accuracy; flag any component whose reward rises
while accuracy stalls (e.g.\ the length term incentivizing short wrong answers).}

% =============================================================================
\section{Analysis and Discussion}
% --- SCAFFOLD ---
\todo{Structure this around the questions the design was built to answer:}
\begin{itemize}
  \item \todo{Which components produced statistically/effect-size-meaningful
  gains, and which did not?}
  \item \todo{pass@1 vs.\ pass@$n$: does the reward shaping expand reasoning or
  sharpen the existing distribution~\citep{yue2025}?}
  \item \todo{The length-reward trade-off (unconditional vs.\ correctness-gated),
  if the sub-experiment was run.}
\end{itemize}

\paragraph{Limitations.}
This study uses \res{3} seeds, a single model size, a single task family
(grade-school and competition math), and a parameter-efficient (QLoRA) rather
than full fine-tuning setup; effect sizes are therefore more reliable than
$p$-values, and replication at larger scale and with full fine-tuning is left to
future work. The partial-credit signal is intentionally simple (number overlap)
and is not a learned process-reward model.

% =============================================================================
\section{Conclusion}
% --- SCAFFOLD ---
We presented a controlled, multi-seed ablation isolating the contribution of
individual reward components in GRPO for small-model mathematical reasoning.
\todo{State the headline finding and its practical implication for reward design
in one or two sentences.}

% =============================================================================
\bibliographystyle{plainnat}
\bibliography{references}

\end{document}
